\documentclass[aspectratio=169]{beamer}

\usepackage[spanish, es-tabla]{babel}
\usepackage{graphicx}
\usepackage{amsmath}
\usepackage{booktabs}
\usepackage{biblatex}
\addbibresource{refs.bib}

\usetheme{Madrid}
\usecolortheme{default}
\renewcommand{\familydefault}{\rmdefault}
\setbeamertemplate{navigation symbols}{}

\title{Simulación Monte Carlo de Casos Nuevos de COVID-19}
\subtitle{Estimación por zona geográfica}
\author{Gustavo A. Calvo Gutiérrez}
\institute{Universidad de Costa Rica}
\date{San José, 2026}

\begin{document}

\begin{frame}
  \titlepage
\end{frame}

\begin{frame}{Pregunta de investigación}
  \begin{block}{Pregunta}
    ¿Cómo puede un modelo de simulación Monte Carlo, aplicado por zona
    geográfica y basado en Xie \cite{Xie2020}, estimar la cantidad esperada de
    casos nuevos de COVID-19 a partir de series temporales observadas?
  \end{block}

  \begin{itemize}
    \item Se trabaja con datos nacionales del Oxford COVID-19 Government
          Response Tracker \cite{Hale2021}.
    \item La simulación resume la incertidumbre mediante promedios y
          percentiles.
  \end{itemize}
\end{frame}

\begin{frame}{Objetivo}
  \begin{block}{Objetivo general}
    Desarrollar un modelo de simulación Monte Carlo para estimar casos nuevos
    de COVID-19 en distintas zonas geográficas.
  \end{block}

  \begin{itemize}
    \item Adaptar el procedimiento de Xie \cite{Xie2020}.
    \item Construir series diarias por país.
    \item Estimar \(R_t\) desde casos observados.
    \item Comparar trayectorias simuladas con casos reales.
  \end{itemize}
\end{frame}

\begin{frame}{Datos utilizados}
  \begin{itemize}
    \item Países: Australia, Brasil, Canadá, China, Reino Unido, India y
          Estados Unidos.
    \item Columnas usadas: \texttt{CountryName}, \texttt{CountryCode},
          \texttt{Jurisdiction}, \texttt{Date} y \texttt{ConfirmedCases}.
    \item Se conservan solo observaciones nacionales:
          \texttt{Jurisdiction = NAT\_TOTAL}.
    \item Los casos nuevos se calculan como diferencias diarias no negativas
          sobre casos acumulados.
  \end{itemize}
\end{frame}

\begin{frame}{Estimación de \(R_t\)}
  \begin{itemize}
    \item Los datos no incluyen \(R_t\) directamente.
    \item Se aproxima con el cociente entre:
      \[
        \frac{\text{promedio móvil de 7 días}}{\text{promedio móvil desplazado 5 días}}
      \]
    \item El valor se limita al intervalo \([0,2]\) para reducir inestabilidad
          numérica.
    \item La serie inicia cuando hay al menos 100 casos nuevos diarios.
  \end{itemize}
\end{frame}

\begin{frame}{Modelo estocástico}
  Para un día \(t\), si hay \(I_t\) casos nuevos simulados:

  \[
    S_t \sim \operatorname{Poisson}(I_t R_t)
  \]

  \begin{itemize}
    \item \(S_t\): contagios secundarios generados desde el día \(t\).
    \item Los retrasos se asignan con una binomial negativa.
    \item El proceso se repite muchas veces para aproximar el valor esperado:
      \[
        E[K_{z,t}] \approx \frac{1}{N}\sum_{k=1}^{N}K_{z,t}^{(k)}.
      \]
  \end{itemize}
\end{frame}

\begin{frame}{Flujo computacional}
  \begin{enumerate}
    \item Preparar series por país.
    \item Estimar casos nuevos diarios y \(R_t\).
    \item Ejecutar 1000 simulaciones por país.
    \item Resumir promedios y percentiles diarios.
    \item Generar figuras comparando simulación y observación.
  \end{enumerate}
\end{frame}

\begin{frame}{Resultados por país}
  \begin{table}
    \centering
    \small
    \begin{tabular}{lrrr}
      \toprule
      País & \(R_t\) medio & Pico mediano & Acumulado mediano \\
      \midrule
      Australia & 1.069 & 237548 & 37015378 \\
      Brasil & 1.059 & 79706 & 18278252 \\
      Canadá & 1.039 & 21632 & 3057878 \\
      China & 1.080 & 632952 & 25822344 \\
      Reino Unido & 1.047 & 92554 & 15543702 \\
      India & 1.032 & 110546 & 15048894 \\
      Estados Unidos & 1.044 & 74788 & 11410956 \\
      \bottomrule
    \end{tabular}
  \end{table}
\end{frame}

\begin{frame}{Ejemplo: Estados Unidos}
  \centering
  \includegraphics[width=0.9\textwidth]{../report/generated/figures/USA_daily_cases.png}
\end{frame}

\begin{frame}{Ejemplo: Brasil}
  \centering
  \includegraphics[width=0.9\textwidth]{../report/generated/figures/BRA_daily_cases.png}
\end{frame}

\begin{frame}{Interpretación}
  \begin{itemize}
    \item El modelo reproduce cualitativamente algunos periodos de crecimiento.
    \item No es una reconstrucción exacta de los casos reales.
    \item La simulación se alimenta de su propia trayectoria después del primer
          día.
    \item Brotes importados, cambios de reporte y heterogeneidad regional no se
          modelan explícitamente.
  \end{itemize}
\end{frame}

\begin{frame}{Conclusiones}
  \begin{itemize}
    \item El enfoque Monte Carlo permite estimar trayectorias esperadas y
          rangos de incertidumbre.
    \item La disponibilidad de variables condiciona la forma final del modelo.
    \item Una versión predictiva debería incorporar reinyección de casos
          observados y calibración por ventanas.
  \end{itemize}
\end{frame}

\begin{frame}[allowframebreaks]{Referencias}
  \printbibliography
\end{frame}

\end{document}
